% Generated draft: requires mathematical and source review.
\documentclass[11pt]{article}
\usepackage[margin=25mm]{geometry}
\usepackage{fontspec}
\setmainfont{Latin Modern Roman}
\usepackage{amsmath,amssymb,mathtools,unicode-math}
\setmathfont{Latin Modern Math}
\usepackage{enumitem,needspace,xurl,hyperref}
\setlength{\parindent}{0pt}
\setlength{\parskip}{5pt}
\setlength{\emergencystretch}{3em}
\newcommand{\N}{\mathbb{N}}
\newcommand{\Z}{\mathbb{Z}}
\newcommand{\Q}{\mathbb{Q}}
\newcommand{\R}{\mathbb{R}}
\newcommand{\C}{\mathbb{C}}
\newcommand{\F}{\mathbb{F}}
\newcommand{\A}{\mathbb{A}}
\newcommand{\PP}{\mathbb{P}}
\newcommand{\E}{\mathbb{E}}
\newcommand{\eps}{\varepsilon}
\newcommand{\dd}{\,\mathrm{d}}
\begin{document}
\section*{UnsolvedMath ID: 30001978}
\textbf{Source problem code:} catalogue entry
\subsection{Definitions and mathematical statement}
Let \(\mathrm{PA}\) denote first-order Peano Arithmetic, with induction for every formula in its language. The usual arithmetic content of Fermat's Last Theorem is
\[
\forall n\in\N\;\forall x,y,z\in\N,\qquad
\bigl(n\geq 3\land x\geq 1\land y\geq 1\land z\geq 1\bigr)
\Longrightarrow x^n+y^n\neq z^n.
\]
Here exponentiation has its usual recursively defined meaning; expressing this statement in the language of \(\mathrm{PA}\) requires an arithmetic encoding of exponentiation.

In a set-theoretic setting, a cardinal \(\kappa\) is \emph{strongly inaccessible} if it is uncountable, regular, and a strong limit: \(\operatorname{cf}(\kappa)=\kappa\) and \(2^\lambda<\kappa\) for every cardinal \(\lambda<\kappa\).

The supplied discussion prompt asks: ``Should number theorists care whether their abstract structures imply strongly inaccessible cardinals, and can logic yield useful number-theoretic results or clarify the strength of Fermat's Last Theorem over Peano Arithmetic?''

This prompt does not specify a mathematical proposition with determined hypotheses and quantifiers. In particular, it names neither the abstract structures nor their governing axioms, specifies no background theory for the claimed implication to inaccessible cardinals, and gives no mathematical criterion for usefulness. The phrase ``strength of Fermat's Last Theorem over Peano Arithmetic'' could concern provability in \(\mathrm{PA}\), equivalence over a specified base theory, or another comparison of axiomatic strength; the supplied text does not select one. The displayed arithmetic statement and the definitions above clarify terminology only and do not supply these missing choices. Consequently, no unique precise conjecture can be recovered from this record.

The supplied statement-verification note classifies the report passage as a historical, pedagogical, or methodological discussion prompt rather than an unresolved mathematical proposition, and marks the statement as unrecoverable.
\subsection{Short English statement}
The record asks broadly whether strong set-theoretic assumptions and logical methods matter to number theory, including the axiomatic strength of Fermat's Last Theorem relative to Peano Arithmetic. It does not specify enough data to determine a precise mathematical claim. Its supplied open label is a curation label, not a verified assertion that a particular conjecture remains unresolved.
\subsection{Sources}
\begin{itemize}
\item Supplied local UnsolvedMath registry record, permanent ID \(30001978\), title: \emph{Logical Strength and Usefulness in Number Theory}; exact source problem code: \texttt{catalogue entry}. The record belongs to the Oberwolfach Reports Open Problems collection and gives the proposed year as 2012; no individual proposer is supplied.
\item \emph{Explicit Versus Tacit Knowledge in Mathematics} (2012), Oberwolfach Reports, pp.~131--210, DOI: \texttt{10.4171/owr/2012/04}. Supplied source URL and literature-source link, labelled ``original Oberwolfach report'': \url{https://doi.org/10.4171/owr/2012/04}. The supplied background describes Oberwolfach Reports as documenting mathematical workshops and their open-problem discussions.
\item Supplied literature assessment, checked exactly \texttt{2026-08-21}: status ``open''; the assessment describes a broad foundational question about strong axioms and formal independence, rather than a single precisely stated theorem, and says that no resolution applies to the extracted text. This is supplied dated curation, not independently verified. The background expressly calls for independent verification of time-sensitive status and cited claims.
\item Supplied statement-verification note: the official history-of-mathematics workshop report treats the passage as a historical, pedagogical, or methodological discussion prompt. No separate clean statement is promoted; the original record is retained for provenance. This characterization is supplied metadata and has not been independently verified here.
\end{itemize}
\end{document}
